取特征线标号 $\psi=G(z)=a(z)-\ell\ln z$。
由于 $G'(z)=(r-z)/z<0$，$G$ 将 $(h,e)$ 光滑地一一映到
$(G(e),G(h))$，其逆映射记为 $z=z(\psi)$。
在开集
\[
\{(s,\psi):s>h,\ G(e)<\psi<G(h),\ \psi+\ell\ln s>0\}
\]
上定义光滑函数
\[
L(s,\psi)=\ln\frac{s-h}{s-r}-\ln(\psi+\ell\ln s),\qquad
F(s,\psi)=L(s,\psi)-L(z(\psi),\psi).
\]
其中 $\psi+\ell\ln z(\psi)=a(z(\psi))>0$，故第二项有定义。
每个非平凡根都满足 $s=\sigma(z)>z>h$ 和
$\psi+\ell\ln s=j(z)>0$，因而位于该开集内；
隐函数定理只在这些根的局部邻域使用，不在 $z=e$ 的合并端点使用。
根方程为 $F(s,\psi)=0$。
